% Compact IEEE-style tables for the networking letter.
% Table II values are measured locally (20 repetitions), not literature values.
% Requires: \usepackage{booktabs}

% Include in Section II-C (Design Goals and Gap Analysis), preferably after contributions.
\begin{table}[t]
\caption{Comparison of the base paper and the proposed design.}
\label{tab:method-comparison}
\centering
\footnotesize
\setlength{\tabcolsep}{3.4pt}
\begin{tabular}{p{2.35cm}ccccc}
\toprule
\textbf{Scheme} & \textbf{Std. KEM} & \textbf{Std. sig.} & \textbf{Formal check} & \textbf{Grouping} & \textbf{Multi-seed stats} \\
\midrule
Base AGS-PBFT~\cite{aslam2025} & ML-KEM precursor & No$^{a}$ & No & Static score & No \\
Recent PQC-V2X schemes~\cite{rasheed2025,mishra2025,asim2026} & Varies & Varies & Varies & Varies & Varies \\
Proposed & ML-KEM-768 & ML-DSA-65 & ProVerif + proof sketch & Online ML + fallback & Paired CI/tests$^{b}$ \\
\bottomrule
\multicolumn{6}{p{7.2cm}}{\scriptsize $^{a}$The source paper names signature tokens but does not specify a verifiable signature construction. $^{b}$Report only after final SUMO/Veins/OMNeT++ runs with $\geq30$ untouched paired seeds.} \\
\end{tabular}
\end{table}

% Include in Section V-C (Formal Security Verification), immediately after the theorem/proof sketch.
\begin{table}[t]
\caption{Security-property mapping and verification target.}
\label{tab:security-mapping}
\centering
\footnotesize
\setlength{\tabcolsep}{4.0pt}
\begin{tabular}{p{2.0cm}p{3.0cm}p{2.25cm}}
\toprule
\textbf{Property} & \textbf{Mechanism} & \textbf{Formal check} \\
\midrule
Mutual authentication & TA-validated ML-DSA credentials; signed \texttt{HELLO}, \texttt{CHALLENGE}, and \texttt{CONFIRM} & Injective correspondence \\
Session-key secrecy & Fresh ML-KEM ephemeral key, transcript hash, and KDF & Secrecy / observational equivalence \\
Replay resistance & Session identifier, nonces, timestamps, and replay cache & Authentication correspondence plus replay-cache rule \\
PBFT integrity & ML-DSA-signed votes; distinct-sender counting; $2f+1$ quorum & Protocol rule under $n\geq3f+1$ \\
\bottomrule
\end{tabular}
\end{table}

% Include in Section VI-A (Implementation and Measurement Setup), before network-level results.
\begin{table}[t]
\caption{Measured PQC overhead on the experiment machine (20 repetitions).}
\label{tab:pqc-overhead}
\centering
\footnotesize
\setlength{\tabcolsep}{4.8pt}
\begin{tabular}{llrr}
\toprule
\textbf{Primitive} & \textbf{Operation} & \textbf{Mean (ms)} & \textbf{Output / key size} \\
\midrule
ML-KEM-768 & Key generation & 10.764 & $pk$: 1184 B; $sk$: 2400 B \\
ML-KEM-768 & Encapsulation & 0.612 & $ct$: 1088 B; $ss$: 32 B \\
ML-KEM-768 & Decapsulation & 0.683 & $ss$: 32 B \\
ML-DSA-65 & Key generation & 1.218 & $pk$: 1952 B; $sk$: 4032 B \\
ML-DSA-65 & Signing & 5.860 & signature: 3309 B \\
ML-DSA-65 & Verification & 1.660 & -- \\
\bottomrule
\multicolumn{4}{p{7.35cm}}{\scriptsize CPU-only measurement on AMD Ryzen 5 5600H (3.30 GHz), with maintained \texttt{liboqs-python}; no GPU or external hardware was used. Values are microbenchmarks, not end-to-end V2X latency.} \\
\end{tabular}
\end{table}
